\chapter{Первый гейт: представление наблюдаемого сектора}

\section{Эмпирическая target-таблица}

До попытки вывода фиксируется target одного поколения. Она служит тестом
representation, а не доказанным следствием S2T:
\begin{center}
\begin{tabular}{cccc}
\toprule
Поле & \(SU(3)_c\) & \(SU(2)_L\) & \(Y\) \\
\midrule
\(Q_L\) & \(\mathbf 3\) & \(\mathbf 2\) & \(1/6\) \\
\(u_R\) & \(\mathbf 3\) & \(\mathbf 1\) & \(2/3\) \\
\(d_R\) & \(\mathbf 3\) & \(\mathbf 1\) & \(-1/3\) \\
\(L_L\) & \(\mathbf 1\) & \(\mathbf 2\) & \(-1/2\) \\
\(e_R\) & \(\mathbf 1\) & \(\mathbf 1\) & \(-1\) \\
\(\nu_R\) & \(\mathbf 1\) & \(\mathbf 1\) & \(0\) \\
\(H\) & \(\mathbf 1\) & \(\mathbf 2\) & \(1/2\) \\
\bottomrule
\end{tabular}
\end{center}

Нулевая гипотеза для первого перебора --- стандартный почти-коммутативный
класс конечных алгебр, содержащий комплексный, кватернионный и цветовой
матричный блоки. Его использование допустимо как baseline. Оно не является
новым открытием и не закрывает задачу происхождения трёх поколений или
Yukawa-матриц.

\section{Representation pass}

Кандидат \((\mathcal A_F,H_F,J_F,\gamma_F)\) проходит первый гейт только
если:
\begin{enumerate}
 \item unitary group после unimodularity/central quotient имеет локальную
 алгебру \(su(3)\oplus su(2)\oplus u(1)\);
 \item left/right modules воспроизводят target-таблицу без post hoc
 переназначения зарядов;
 \item внутренние флуктуации создают один Higgs doublet либо заранее
 объявленное расширение;
 \item допустимые рёбра \(D_F\) дают gauge-invariant Yukawa vertices;
 \item real structure создаёт antiparticle sectors без удвоения физических
 степеней свободы;
 \item sterile neutrino и Majorana edge имеют однозначно объявленный статус.
\end{enumerate}

\section{Немедленные anomaly tests}

Для каждого поколения должны обращаться в нуль коэффициенты
\begin{equation}
 \mathcal A_{333},\quad
 \mathcal A_{221},\quad
 \mathcal A_{331},\quad
 \mathcal A_{111},\quad
 \mathcal A_{\mathrm{grav}^2 1},
\end{equation}
с учётом того, что правые фермионы при подсчёте chiral anomaly записываются
как левые charge-conjugate states. Проверка проводится из выведенного
ledger, а не подстановкой заранее известной таблицы после построения.

\begin{protocol}[Первый вычислительный пакет]
Следующий машинный этап должен:
\begin{enumerate}
 \item перечислить минимальные finite-algebra candidates;
 \item построить все irreducible left/right bimodules в заданной границе
 размерности;
 \item вычислить gauge algebra, hypercharge generator и anomaly polynomial;
 \item отфильтровать кандидаты по target-таблице;
 \item отдельно отметить, какие свойства были выведены, а какие внесены как
 empirical constraints.
\end{enumerate}
\end{protocol}

\begin{nogocriterion}[Baseline не является предсказанием]
Если target-таблица целиком вставлена в \(H_F\), а затем из неё
восстановлено действие Стандартной модели, получен корректный baseline, но
не объяснение наблюдаемого спектра. Для статуса нового результата требуется
либо уникальность представления в заранее заданном классе, либо новое
ограничение на нормировки, массы, mixing или neutrino sector, прошедшее
blind-test.
\end{nogocriterion}

\section{Ближайшая цель}

Первой реальной вехой Version IV.SM является не масса частицы, а
машинно воспроизводимый classification result:
\begin{equation}
 \{\text{допустимые finite geometries}\}
 \longrightarrow
 \{\text{anomaly-free observed candidates}\}.
\end{equation}
Только после этого разрешён расчёт spectral coupling relations и RG.